% Einheit 3 von 3 – Faktorisieren (bank/binomische-formeln/e3.jsonl, zusammenbau v0.1)
%% TODO Zweigzeile Teil 1: was hier gelernt wird (Satz in Schülersprache aus dem Einheitstitel) – Einheit 3
\einheitenkopf[e3]{Einheit 3 von 3 · Faktorisieren}
\zweigzeile{ab Kl. 8, je nach Buch bis Kl. 10 · keine P10-Aufgabe}
\setcounter{aufgabe}{16}

%% TODO Titel: Ich-kann-Satz fehlt in der Bank (Platzhalter: Kettenname) – Nr. 17
%% TODO Anweisung: ein Satz über der ersten Teilaufgabe fehlt in der Bank – Nr. 17
\begin{aufgabe}{Faktorisieren}
\begin{teile}
\teil $x^2 + 16x + 64$ – Quadrate mit Wurzeln? Passt das Mittelglied? \janein Wurzeln \leerfeld und \leerfeld
\end{teile}
%% TODO Darstellung für schwach fehlt in der Bank (binomische-formeln-e3-k1-s1-v1; Abschnitt „Für schwache Schüler“)
\swz{$\text{Faktorisiere: $x^2 - 36$}$}{}
%% TODO Darstellung für schwach fehlt in der Bank (binomische-formeln-e3-k1-s1-v2; Abschnitt „Für schwache Schüler“)
\swz{$\text{Faktorisiere: $x^2 - 81$}$}{}
%% TODO Darstellung für schwach fehlt in der Bank (binomische-formeln-e3-k1-s1-v3; Abschnitt „Für schwache Schüler“)
\swz{$\text{Faktorisiere: $x^2 - 16$}$}{}
%% TODO Darstellung für schwach fehlt in der Bank (binomische-formeln-e3-k1-s1-v4; Abschnitt „Für schwache Schüler“)
\swz{$\text{Faktorisiere: $x^2 - 1$}$}{}
%% TODO Darstellung für schwach fehlt in der Bank (binomische-formeln-e3-k1-s1-v5; Abschnitt „Für schwache Schüler“)
\swz{$\text{Faktorisiere: $x^2 - 121$}$}{}
\swfrage{Was bleibt gleich, was ändert sich?}
%% TODO Darstellung für schwach fehlt in der Bank (binomische-formeln-e3-k1-s2-v1; Abschnitt „Für schwache Schüler“)
\swz{$\text{Faktorisiere: $x^2 + 8x + 16$}$}{}
%% TODO Darstellung für schwach fehlt in der Bank (binomische-formeln-e3-k1-s3-v1; Abschnitt „Für schwache Schüler“)
\swz{$\text{Faktorisiere: $x^2 - 18x + 81$}$}{}
%% TODO Darstellung für schwach fehlt in der Bank (binomische-formeln-e3-k1-s4-v1; Abschnitt „Für schwache Schüler“)
\swz{$\text{Faktorisiere: $16x^2 - 9$}$}{}
%% TODO Darstellung für schwach fehlt in der Bank (binomische-formeln-e3-k1-s5-v1; Abschnitt „Für schwache Schüler“)
\swz{$\text{Faktorisiere: $2x^2 - 32$}$}{}
%% TODO Darstellung für schwach fehlt in der Bank (binomische-formeln-e3-k1-s6-v1; Abschnitt „Für schwache Schüler“)
\swz[3]{Lässt sich $x^2 + 64$ mit einer binomischen Formel faktorisieren? Begründe.}{}
%% TODO Darstellung für schwach fehlt in der Bank (binomische-formeln-e3-k1-s7-v1; Abschnitt „Für schwache Schüler“)
\swz[3]{Stimmt $x^2 - 16x + 64 = (x - 8)^2$? Mach die Probe durch Ausmultiplizieren.}{}
%% TODO Darstellung für schwach fehlt in der Bank (binomische-formeln-e3-k1-s8-v1; Abschnitt „Für schwache Schüler“)
\swz{$\text{Löse durch Faktorisieren: $x^2 - 144 = 0$}$}{}
%% TODO Darstellung für schwach fehlt in der Bank (binomische-formeln-e3-k1-s9-v1; Abschnitt „Für schwache Schüler“)
\swz{$\text{Schreib mit quadratischer Ergänzung als Quadrat plus Zahl: $x^2 + 8x + 3$}$}{}
\swz[3]{$\text{Faktorisiere vollständig und mach die Probe: $3x^2 - 48$}$}{}
\end{aufgabe}

%% TODO Titel: Ich-kann-Satz fehlt in der Bank (Platzhalter: Kettenname) – Nr. 18
%% TODO Anweisung: ein Satz über der ersten Teilaufgabe fehlt in der Bank – Nr. 18
\begin{aufgabe}{Faktorisieren – Fehler finden, Begründen}
\swz[3]{Finde den Fehler und rechne richtig: \rechnung{x^2 + 1 &= (x + 1)^2}}{}
\swfrage{Warum lässt sich $x^2 + 16$ nicht mit einer binomischen Formel faktorisieren? Begründe.}
\end{aufgabe}

\uebersichtskasten{Faktorisieren: eine Summe in ein Produkt verwandeln – erst gemeinsamen Faktor ausklammern, dann nach einer binomischen Formel suchen. \\ Drei Glieder, zwei Quadrate, Mittelglied ist das doppelte Produkt der Wurzeln: x² + 10x + 25 = (x + 5)² \quad Minus im Mittelglied: x² − 10x + 25 = (x − 5)² \\ Zwei Quadrate mit Minus dazwischen: x² − 25 = (x + 5)·(x − 5) \\ Erst ausklammern: 3x² − 75 = 3·(x² − 25) = 3·(x + 5)·(x − 5) \\ Probe: ausmultiplizieren. Passt das Mittelglied nicht (x² + 7x + 25), gibt es keine binomische Formel.}
